\documentclass[10pt,a4paper]{amsart}
\usepackage{amsmath,amssymb,amsthm,graphicx}
\usepackage{verbatim}
\newenvironment{reference}{\comment}{\endcomment}
%\newenvironment{reference}{}{}
\newenvironment{slogan}{\comment}{\endcomment}
\newenvironment{history}{\comment}{\endcomment}

% For commutative diagrams we use Xy-pic
\usepackage[all]{xy}

% We use 2cell for 2-commutative diagrams.
\xyoption{2cell}
% Theorem environments.
%
\theoremstyle{plain}
\newtheorem{theorem}[subsection]{Theorem}
\newtheorem{proposition}[subsection]{Proposition}
\newtheorem{lemma}[subsection]{Lemma}

\theoremstyle{definition}
\newtheorem{definition}[subsection]{Definition}
\newtheorem{example}[subsection]{Example}
\newtheorem{exercise}[subsection]{Exercise}
\newtheorem{situation}[subsection]{Situation}

\theoremstyle{remark}
\newtheorem{remark}[subsection]{Remark}
\newtheorem{remarks}[subsection]{Remarks}

\numberwithin{equation}{subsection}

% Macros
%
\def\lim{\mathop{\mathrm{lim}}\nolimits}
\def\colim{\mathop{\mathrm{colim}}\nolimits}
\def\Spec{\mathop{\mathrm{Spec}}}
\def\Hom{\mathop{\mathrm{Hom}}\nolimits}
\def\Ext{\mathop{\mathrm{Ext}}\nolimits}
\def\SheafHom{\mathop{\mathcal{H}\!\mathit{om}}\nolimits}
\def\SheafExt{\mathop{\mathcal{E}\!\mathit{xt}}\nolimits}
\def\Sch{\mathit{Sch}}
\def\Mor{\mathop{\mathrm{Mor}}\nolimits}
\def\Ob{\mathop{\mathrm{Ob}}\nolimits}
\def\Sh{\mathop{\mathit{Sh}}\nolimits}
\def\NL{\mathop{N\!L}\nolimits}
\def\CH{\mathop{\mathrm{CH}}\nolimits}
\def\proetale{{pro\text{-}\acute{e}tale}}
\def\etale{{\acute{e}tale}}
\def\QCoh{\mathit{QCoh}}
\def\Ker{\mathop{\mathrm{Ker}}}
\def\Im{\mathop{\mathrm{Im}}}
\def\Coker{\mathop{\mathrm{Coker}}}
\def\Coim{\mathop{\mathrm{Coim}}}

% Boxtimes
%
\DeclareMathSymbol{\boxtimes}{\mathbin}{AMSa}{"02}

%
% Macros for moduli stacks/spaces
%
\def\QCohstack{\mathcal{QC}\!\mathit{oh}}
\def\Cohstack{\mathcal{C}\!\mathit{oh}}
\def\Spacesstack{\mathcal{S}\!\mathit{paces}}
\def\Quotfunctor{\mathrm{Quot}}
\def\Hilbfunctor{\mathrm{Hilb}}
\def\Curvesstack{\mathcal{C}\!\mathit{urves}}
\def\Polarizedstack{\mathcal{P}\!\mathit{olarized}}
\def\Complexesstack{\mathcal{C}\!\mathit{omplexes}}
% \Pic is the operator that assigns to X its picard group, usage \Pic(X)
% \Picardstack_{X/B} denotes the Picard stack of X over B
% \Picardfunctor_{X/B} denotes the Picard functor of X over B
\def\Pic{\mathop{\mathrm{Pic}}\nolimits}
\def\Picardstack{\mathcal{P}\!\mathit{ic}}
\def\Picardfunctor{\mathrm{Pic}}
\def\Deformationcategory{\mathcal{D}\!\mathit{ef}}
\title{An essentially finitely presented local homomorphism is a colimit of essentially finite type ones with localised transition maps}
\author{The Stacks Project Authors}
\date{stacks-project:algebra.tex:local-snapshot-20261012; GNU FDL 1.2 or later}
\begin{document}
\maketitle
\noindent\textit{Editorial note.} All delivered mathematical text is the author's own, copied from the pinned Stacks Project chapter file algebra.tex; no statement or proof step was rewritten, reordered or removed.\par\medskip
\noindent\textit{Editorial note.} Preamble provenance, as required for this source: the frozen snapshot stages the chapter's own preamble file beside it (preamble.tex, 217 lines, 8116 bytes, the same shared Stacks-wide preamble the published chow and etale-cohomology cards use); the first line of the pinned chapter file is \textbackslash{}input\{preamble\}, so no preamble text lives inside the chapter itself. The card takes that file's lines 8-18 (the verbatim package and the three \textbackslash{}comment-based environments reference/slogan/history at lines 9-12, section 13.2, together with the Xy-pic setup) and its lines 154 to the runtime end of the file (the native theorem declarations and mathematical macros). Not one character was changed, no package was added, no macro was defined, and the excerpt serves only as preamble, never as a span coordinate.\par\medskip
\noindent\textit{Editorial note.} The same-chapter pointers are located at their \textbackslash{}label lines inside the pinned chapter file; the cross-chapter pointers are located in the project's own tag list staged in the snapshot (sources/stacks-chow/tags.txt), at the verbatim lines carrying their canonical tags.\par\medskip
\noindent\textit{Editorial note.} The counted proof ends by leaving the verification of the four listed properties to the reader; that sentence is kept verbatim inside the span and the construction it applies to is written out in full.\par\medskip
\noindent\textit{Editorial note.} The counted spans contain no \textbackslash{}ref and no \textbackslash{}cite, so the card carries no navigation stub and no bibliography.\par\medskip
\noindent\textit{Allowed standard prerequisites (editorial).}\begin{itemize}
\item Essentially of finite presentation as a localisation of a finitely presented algebra. Used at the start of the counted proof (lines 32585-32590). Supplies the explicit presentation and the prime at which it is localised.
\item Descending finitely many generators along a colimit presentation of the base. Used in the middle of the counted proof (lines 32592-32599). Moves the finitely many coefficients of the relations down to a single index of the colimit.
\end{itemize}\medskip
\begin{lemma}
\label{lemma-limit-essentially-finite-presentation}
Suppose $R \to S$ is a local homomorphism of local rings.
Assume that $S$ is essentially of finite presentation over $R$.
Then there exists a directed set $(\Lambda, \leq)$, and
a system of local homomorphism $R_\lambda \to S_\lambda$
of local rings such that
\begin{enumerate}
\item The colimit of the system $R_\lambda \to S_\lambda$
is equal to $R \to S$.
\item Each $R_\lambda$ is essentially of finite type
over $\mathbf{Z}$.
\item Each $S_\lambda$ is essentially of finite type
over $R_\lambda$.
\item For each $\lambda \leq \mu$ the map
$S_\lambda \otimes_{R_\lambda} R_\mu \to S_\mu$
presents $S_\mu$ as the localization of
$S_\lambda \otimes_{R_\lambda} R_\mu$
at a prime ideal.
\end{enumerate}
\end{lemma}

\begin{proof}
By assumption we may choose an isomorphism
$\Phi : (R[x_1, \ldots, x_n]/I)_{\mathfrak q} \to S$
where $I \subset R[x_1, \ldots, x_n]$ is a finitely generated ideal,
and $\mathfrak q \subset R[x_1, \ldots, x_n]/I$ is a prime.
(Note that $R \cap \mathfrak q$
is equal to the maximal ideal $\mathfrak m$ of $R$.)
We also choose generators $f_1, \ldots, f_m \in I$ for the ideal $I$.
Write $R$ in any way as a colimit $R = \colim R_\lambda$
over a directed set $(\Lambda, \leq )$, with each $R_\lambda$
local and essentially of finite type over $\mathbf{Z}$.
There exists some $\lambda_0 \in \Lambda$ such that $f_j$ is the image
of some $f_{j, \lambda_0} \in R_{\lambda_0}[x_1, \ldots, x_n]$.
For all $\lambda \geq \lambda_0$ denote
$f_{j, \lambda} \in R_{\lambda}[x_1, \ldots, x_n]$ the image
of $f_{j, \lambda_0}$. Thus we obtain a system of ring maps
$$
R_\lambda[x_1, \ldots, x_n]/(f_{1, \lambda}, \ldots, f_{m, \lambda})
\to
R[x_1, \ldots, x_n]/(f_1, \ldots, f_m) \to S
$$
Set $\mathfrak q_\lambda$ the inverse image of $\mathfrak q$.
Set $S_\lambda = (R_\lambda[x_1, \ldots, x_n]/
(f_{1, \lambda}, \ldots, f_{m, \lambda}))_{\mathfrak q_\lambda}$.
We leave it to the reader to see that this works.
\end{proof}

\end{document}
